\section*{Sesquilinear forms:}

Let A be a ring, \(\xi \to \overline{\xi}\) an involutive antiautomorphism of A, that is, a bijection of A onto itself such that \(\overline{(\xi + \eta)} = \overline{\xi} + \overline{\eta}\) , \(\overline{(\xi\eta)} = \overline{\eta} . \overline{\xi}\) , \(\overline{\overline{\xi}} = \xi\) for all \(\xi, \eta\) in A. A sesquilinear form \(\Phi\) on a left A-module E is a map from E × E into A such that

\[
\begin{array}{c c} \Phi (x + x ^ {\prime}, y) = \Phi (x, y) + \Phi (x ^ {\prime}, y), & \Phi (x, y + y ^ {\prime}) = \Phi (x, y) + \Phi (x, y ^ {\prime}) \\ \Phi (\alpha x, y) = \alpha \Phi (x, y), & \Phi (x, \alpha y) = \Phi (x, y) \bar {\alpha} \end{array}
\]

for all \(\alpha\in\mathbf{A},x,x^{\prime},y,y^{\prime}\) in E.

Let \(\varepsilon\) be an element of the center of A. One says that a sesquilinear form \(\Phi\) on E is \(\varepsilon\)-hermitian if \(\Phi(y, x) = \varepsilon \overline{\Phi(x, y)}\) for all \(x \in E\) , \(y \in E\) ; if \(\varepsilon = 1\) (resp. \(\varepsilon = -1\) ) one says that \(\Phi\) is hermitian (resp. anti-hermitian).

When the antiautomorphism \(\xi\to\overline{\xi}\) is the identity (which implies that the ring A is commutative), the corresponding sesquilinear forms are the bilinear forms. One then says "symmetric" instead of "hermitian", and "antisymmetric" instead of "antihermitian" (cf.~Chap.~III). A bilinear form \(\Phi\) such that \(\Phi(x,x)=0\) is said to be alternating; it is then antisymmetric, and the converse is true when A is a field of characteristic \(\neq2\) (cf.~Chap.~III).

One says that a form \(\varepsilon\)-hermitian satisfies condition (T) if for all \(x \in E\) , there exists \(\lambda \in A\) such that \(\Phi(x, x) = \lambda + \varepsilon \overline{\lambda}\) . This condition is always satisfied when \(\varepsilon = 1\) and \(A\) is a field of characteristic \(\neq 2\) , or when \(\Phi\) is alternating.

